\chapter{Gender-fair language in (machine) translation} \label{Chapter2-3}

The visibility of non-binary individuals has increased in the last few years. Celebrities such as Sam Smith and Demi Lovato came out as non-binary and television series such as \textit{Sex Education} or \textit{One Day at a Time} are more frequently featuring non-binary characters. In some languages, this represents no challenge: for example, in English where singular \textit{they} can be used as a gender-neutral pronoun (\cite{apasingularthey}) and most nouns are not gender-specific. In others, gender must be marked and new GFL strategies have been proposed, such as the inclusive \textit{e} in Spanish (\cite{lopez2019tu}). The chapter delivers an introduction to the topic and presents GFL strategies in three different languages, namely English, German, and Italian. It then moves on to present research on GFL in TS and MT.

\section{Introducing gender-fair language} \label{inclusivelanguage}
%Lopez' distinction between direct and indirect non-binary language
\textcite{sczesny2016can} use the term gender-fair and highlight how other expressions such as gender-neutral and gender-inclusive are also used, whereby the two have actually different meanings. Gender-neutral means that linguistic forms are not marked for gender, whereas gender-inclusive refers to linguistic forms that represent all gender identities (\cite[28]{hornscheidt2021wie}). This also reflects the process some languages are undergoing. In English, for example, as a notional gender language that is moderately inflected, speakers are headed towards a degendering, i.e. gender neutrality, to counteract sexism \citep[17]{weatherall2002gender}. For this reason, gender suffixes such as \textit{-ess} in \textit{actress} are becoming obsolete \citep[220]{eckert2013language}. In grammatical gender languages, such as German and Italian, which are morphologically richer, new suffixes to include all gender identities have been introduced (see for example \cite{hornscheidt2021wie} and \cite{vitiello2021linguaggio} for some overviews), which corresponds to a process of engendering. This means that strategies also differ based on the gender system of a language. As a consequence, gender-fair is used in this work as a generic term to include every strategy to foster linguistic equality between the sexes/genders. Note also that gender-inclusive and gender-neutral are often erroneously used interchangeably even though they refer to two distinct processes. 

Similarly, in their work \textcite{lopez2019tu} distinguishes between Direct Non-Binary Language (DNL) and Indirect Non-Binary Language (INL). The former strategy describes those approaches that aim to address non-binary individuals through new gendered suffixes (e.g. the use of \textit{e} and \textit{es} in Spanish), typographical characters (e.g. @), neopronouns (e.g. \textit{elle} and \textit{elles}) and/or neologisms. The latter refers to the use of sentences structured in order to avoid gender markers, for instance through passive constructions and indefinite pronouns. 

An interesting reflection on the term inclusion was offered by \textcite[121]{acanfora2021in}. Its definition as ``the act of including someone or something as part of a group, list, etc., or a person or thing that is included'' \citep{cambridgeinclusion} reveals that there is a power imbalance between those who include and those who are included. It is the former who decide whether to include the latter or not. \textcite[125f]{acanfora2021in} thus proposes a re-branding of this concept by introducing the expression \textit{convivenza delle differenze} (coexistence of differences): \textit{coexistence} suggests that relationships between people are based on mutual respect and equality, while \textit{differences} is used to move away from the concept of norm which is often used as a standard to define what is considered as not normal and therefore inferior. In Italian, for example, other terms to refer to inclusive language, in general, are \textit{linguaggio aperto} and \textit{linguaggio ampio} (lit. open and broad language to illustrate the inclusion of all oppressed identities) and describe a language that is free from stereotypes and respectful of all differences, including gender, race, age, class, and so on (e.g. \cite{gheno2022questione}).

Finally, Subsection \ref{Sexism in language} described how sexist language use covers a different range of phenomena such as the use of masculine generics, asymmetrical forms, or gender stereotypes. The present work focuses on the representation of non-binary individuals in language, who are often excluded from public discourse or misgendered. Expressions like gender-fair and gender-inclusive are consequently to be understood as making them visible and delivering the appropriate linguistic means to refer to them. Though discrimination can be perpetrated on other, subtler levels, e.g. through stereotypes, this is not the focus of the present work. 

%--------------------------------------
%	SUBSECTION 1
%-------------------------------------
\subsection{Gender-fair language in English}
In Subsection \ref{sub:language/gender}, the gender system of the English language was briefly introduced, highlighting how gender is mainly marked in third-person singular pronouns (he/she/it) and in some specific noun pairs (boy/girl) or job titles (actor/actress). In order to represent non-binary individuals, the process English is undergoing is primarily degendering or neutralisation (\cite[17]{weatherall2002gender}, \cite{sczesny2016can}), i.e. markers for gender are being removed and/or gendered words are being replaced by gender-neutral alternatives. This tendency, however, started with second-wave feminism movements and, especially, in the United States where feminists focused on two issues in language use, i.e. the use of masculine generics and of gender-specific titles in contexts where gender is not relevant \citep[73]{kramer2016feminist}.

According to \textcite{apagenderenglish} and \textcite{europarlenglish},\footnote{Note that the list provided by \textcite{europarlenglish} is still based on a binary conception of gender. For a more detailed list of gendered terms with alternative neutral versions, see \textcite{nonbinarywikienglish}.} binary terms and gendered occupational titles must be avoided and replaced by gender-neutral forms, e.g. \textit{housewife} by \textit{homemaker}. More examples can be found in Table \ref{tab:genderenglish} as well as \textcite{apagenderenglish}. Masculine generics should be avoided through gender-neutral words and/or expressions when talking about humans in general, e.g. \textit{human rights} instead of \textit{men's rights} \citep{apagenderenglish} as in Table \ref{tab:malegenericsenglish}. Singular \textit{they} should be applied when talking about a person whose gender is unknown or a hypothetical person whose gender is irrelevant, e.g. ``The client is usually the best judge of the value of their counseling'' \citep{apagenderenglish}. 

\begin{table}
    \caption[Gender-fair alternatives to binary terms in English]{Gender-fair alternatives to binary terms and gendered professions in English \citep{apagenderenglish,europarlenglish}}
    \label{tab:genderenglish}
    \begin{tabular}{ll}
    \lsptoprule
        {Problematic terms}  &  {Preferred terms}  \\
\midrule
        mother, father & parent \\ 
        sister, brother & sibling \\ 
        steward, stewardess & flight attendant \\ 
        waiter, waitress & servant \\ 
        actor, actress & actor \\ 
        %author, authoress & author \\
        chairman & chairperson, chair \\
\lspbottomrule
    \end{tabular}
\end{table}


\begin{table}
    \caption[Gender-fair alternatives to masculine generics in English]{Gender-fair alternatives to masculine generics in English \citep{apagenderenglish,europarlenglish}} 
    \label{tab:malegenericsenglish}
    \begin{tabular}{ll}
    \lsptoprule
       \textbf{Problematic terms}  & \textbf{Preferred terms}  \\
\midrule
        men's rights & human rights \\  
        men, mankind & people, individuals, humanity \\ 
        manpower & workforce \\ 
        statesmen & political leaders \\
\lspbottomrule
    \end{tabular}
\end{table}


More generally, terms and expressions that imply binaries, such as \textit{opposite sex} or \textit{opposite gender} should be avoided. For instance, expressions such as mixed \textit{sex/gender} should be preferred to describe people of different sex/genders in a relationship. Gendered expressions should also be avoided since they convey stereotypes. For instance, \textit{mothering} can be replaced by \textit{parenting} since the first term implies that only mothers take care of their children and that there are only mothers and fathers. For the same reason, \textit{mother tongue} is changed to first \textit{language}. Furthermore, gender should not be specified unless it is relevant to the context of the conversation (see \cite{apagenderenglish}). Likewise, it is better to avoid stereotypical expressions suggesting how sexes/genders are and/or should be. 

An alternative strategy to this process of neutralisation is the introduction of the sometimes so-called neopronouns, a term that is not actually correct as they date back centuries (see \cite{baron2020s} for a detailed historical account). Since the 1780s, there have been many attempts to introduce gender-neutral pronouns in the English language with a uniquely intense ongoing debate (see \cite{baron2020s}, who gathered a wide collection of sources to account for this debate and the creation of more than 200 pronouns). These attempts have their origin in the distinction among genders in the third-person singular pronoun only, while other pronouns and most nouns are gender-neutral. When talking about people in general or a hypothetical person in sentences like ``everyone forgets his passwords'', the gender of the indefinite pronoun is specified through the third-person singular masculine pronoun. The use of the supposedly generic \textit{he} has been deemed as inclusive of women in a time where non-binary people were practically ignored and yet used to exclude women from access to basic rights, such as to vote \citep{baron2020s}. As a consequence, different groups of people started to work without particular success on gender-neutral pronouns to overcome this issue and nowadays some are used specifically by non-binary people (see an overview in Table \ref{tab:pronouns}). 

\begin{table}
    \caption[English non-binary pronouns]{Neopronouns in English (\cite{englishpronouns})}
    \label{tab:pronouns}
    \begin{tabular}{lllll}
    \lsptoprule
         1 & 2 & 3 & 4 & 5 \\
         \midrule
         (f)ae & (f)aer & (f)aer & (f)aers & (f)aersel  \\ 
         e/ey & em & eir & eirs & eirself \\ 
         per & per & pers & pers & perself \\ 
         they & them & their & theirs & themself \\ 
         ve & ver & vis & vis & verself \\ 
         xe & xem & xyr & xyrs & xemself \\ 
         ze/zie & hir & hir & hirs & hirself \\
\lspbottomrule
    \end{tabular}
\end{table}

Grammarians such as \textcite{fisher1753grammar} have historically recommended generic \textit{he} based on \textcite{Lily1549grammar}. The scholar, building on Latin grammar, believed the masculine gender to be worthier than the feminine and the neuter. Note that singular \textit{they} has nevertheless been used since 1375 as a gender-neutral pronoun to refer to a hypothetical person whose gender is irrelevant or unknown \citep{baron2020s}. It is only recently, also thanks to the increasing visibility of non-binary people, that the use of this latter pronoun has been endorsed by language authorities (see for example \cite{apasingularthey} and \cite{merriamsingularthey}) and also became Merriam-Webster's word of the year in 2019 \citep{merriamwordoftheyear}. According to a survey conducted by \textcite{gendercensus2021} on language use by people whose gender is not binary, nearly 80\% of respondents used the pronoun \textit{they} (see Figure \ref{fig:censuspronouns})

%---------IMAGE begins here----------
\begin{figure}[ht]
\centering 
\includegraphics[width=0.8\textwidth]{figures/censuspronouns.pdf}
\caption[Use of gender-neutral pronouns by trans or non-binary people]{Use of gender-neutral pronouns by trans or non-binary people (\cite{gendercensus2021}). Note that respondents could choose more than one pronoun, so percentages add up to more than 100.}
\label{fig:censuspronouns}
\end{figure}
%------------IMAGE ends here------------

The use of forms of address has changed as well. First, \textit{Ms} has been introduced as the feminine equivalent to \textit{Mr} to overcome the sexist distinction between \textit{Mrs} and \textit{Miss} (see Subsection \ref{Sexism in language}) \citep[21]{weatherall2002gender}. Second, a gender-neutral version, i.e. \textit{Mx} has not only been introduced, but also officially recognised by public institutions in the UK (\cite{henry2015mx}). This shows how the English language is undergoing a process of neutralisation to remove gender markers from every word class except pronouns. In this last case, the use of singular \textit{they} is becoming the norm when referring to people of unknown gender, mixed-gender groups and non-binary individuals. 


%--------------------------------------
%	SUBSECTION 2
%-------------------------------------
\subsection{Gender-fair language in German} \label{sec:German}

As a grammatical gender language and hence described as morphologically rich, German is currently undergoing the process of including and representing non-binary people, i.e. multigendering or engendering (\cite[17]{weatherall2002gender}, \cite{sczesny2016can}). Nevertheless, some strategies aim to achieve gender neutrality. Importantly, due to its structure, gender markers are not being removed but rather substituted with typographical characters, such as asterisk, or new suffixes that are meant to be neutral (see \cite[294]{hornscheidt2012feministische}, \cite[53]{hornscheidt2021wie}). This shows, as previously mentioned, the lack of consistency in GFL terminology. For the German language, four main approaches can be identified:\footnote{This classification is based on the commonalities the strategies share across all approaches.}
\begin{enumerate}
\setlength\itemsep{-0.5em}
    \item Gender-neutral rewording: these strategies aim to use the available inventory of language without introducing new pronouns or suffixes, by rephrasing/rewording in gender-neutral ways.
    \item Typographical gender-inclusive strategies: typographical characters such as gender star (*) are introduced to separate masculine from feminine forms of words and thus represent all genders.
    \item Gender-neutral strategies: typographical characters or new suffixes that are meant to be gender-neutral and used in contexts where gender is unknown or irrelevant.
    \item Gender-fair neosystems: introduction of a fourth gender in addition to masculine, feminine, and neuter, and a new inventory of suffixes for (almost) each word class to make non-binary individuals visible and/or to refer to mixed groups of people. 
\end{enumerate}
Other aspects of language change are: (i) pronouns; (ii) forms of address; (iii) binary terms; and (iv) gender stereotypes. They are covered in more detail below. Due to the richness of the debate on gender-inclusivity, it is nevertheless impossible to cover all proposed strategies. The author decided to cover the most used and those found in the literature. However, in social media and more generally on the internet, one can find many more approaches to GFL, see for example \textcite{nibispace}.

The simplest possible approach to GFL is gender-neutral rewording. Sentences can be built such that gender is not marked (\cite[6]{unileipzig2020}, \cite[26]{boka2021richtlinie}), for instance:
\begin{itemize}
\setlength\itemsep{-0.5em}
    \item Collective and singular nouns may be built with specific suffixes and compounds such as -kraft, -\textit{ung}, -\textit{hilfe}, -\textit{person}, -\textit{berechtigte} as in \textit{Lehrkraft} (teacher) or \textit{Fachperson} (expert).
    \item Gender-neutral nouns and nouns referring to functions and institutions may be used, for example \textit{Person}, \textit{Mitglied} (member), \textit{Mensch} (person), \textit{Ministerium} (ministery), \textit{Leitung} (direction), and \textit{Team}.
    \item Plural nouns may be built by nominalising of the present participle, for example \textit{Studierende} and \textit{Lehrende} -- note that such forms have become quite widespread but might not be considered as neutral or inclusive anymore as they are associated with maleness \citep[49]{hornscheidt2021wie}.
    \item Reference to gender may also be avoided by using infinitive and passive constructions.
\end{itemize}

The second approach is the use of typographical characters (see Table \ref{tab:typograpghicalinnovations_german}) to separate masculine from feminine forms of words and make all genders visible. These typographical characters are pronounced as a brief pause in the spoken language. Since there are no official rules on how to build these forms, some creativity is required (\cite[4]{unileipzig2020}, \cite[27]{boka2021richtlinie}). These strategies have now become quite common in written language use (see \cite{boka2021richtlinie}).

\begin{table}[ht]
    \caption[Strategies for gender-inclusive German]{Strategies for gender-inclusive German (adapted from \cite[16]{sprachhandeln2015tun})}
    \label{tab:typograpghicalinnovations_german}
    \noindent\makebox[\textwidth]{
    \resizebox{\textwidth}{!}{\begin{tabular}{cccccc}
    \lsptoprule
         \textbf{Strategy} & \textbf{Nouns} & \textbf{Personal  Pronouns} & \textbf{Possessives} & \textbf{Definite  Articles} & \textbf{Interrogatives}  \\
\midrule
         * & Student*in, Student*innen & sie*er, si*er & ihre*seine & die*der, di*er & welche*r? \\
         \_\footnotemark & Student\_in, Student\_innen & si\_e, er\_sie & ihre\_seine & die\_der, di\_er & welche\_r? \\
         : & Student:in, Student:innen & si:er, er:sie & ihre:seine & die:der, di:er & welche:r? \\
         / & Student/in, Student/innen & si/er, er/sie & ihre/seine & die/der, di/er & welche/r? \\
         ' & Student’in, Student’innen & si’er, er’sie & ihre’seine & die’der, di’er & welche'r? \\
\lspbottomrule
    \end{tabular}}}
\end{table}
\footnotetext{The underscore is also used as \textit{dynamischer Unterstrich} (German for dynamic underscore), i.e. it is placed anywhere within the word to allow a more critical questioning of the binary gender system \citep[304]{hornscheidt2012feministische}.}

In contexts where sex/gender is or should be irrelevant, usual masculine and feminine suffixes can be replaced by typographical characters (\cite[294]{hornscheidt2012feministische}) or new gender-neutral suffixes (\cite[53]{hornscheidt2021wie}; see Table \ref{tab:genderneutrality_german} for an overview of some). The former solutions are attached to the verb stem, and the latter to the noun stem to build nouns. Moreover, they can stand alone to replace some elements, such as pronouns. These strategies are normally used to question the gender binary as well as to irritate and provoke readers (see \cite{hornscheidt2012feministische} for a detailed explanation of how and when to use them). In particular, \textcite[53]{hornscheidt2021wie} propose the use of gender-neutral Ens-forms which derive from the German word M\underline{ens}ch (person, human). Ens-forms are easy to pronounce and provide a complete inventory of suffixes for each word class, hence they might also be classified as a neosystem (see \cite{paolucci-etal-2023-gender}).


\begin{table}[ht!]
    \caption[Selected strategies for gender-neutral German]{Selected strategies for gender-neutral German (\cite[294]{hornscheidt2012feministische}, \cite{sprachhandeln2015tun}, \cite[53]{hornscheidt2021wie})}
    \label{tab:genderneutrality_german}
    \resizebox{\textwidth}{!}{
    \begin{tabular}{cccccc}
    \lsptoprule
        \textbf{Strategy} & \textbf{Nouns} & \textbf{Personal  Pronouns} & \textbf{Possessives} & \textbf{Articles} & \textbf{Indefinite  Pronouns} \\
\midrule
        x & Studierx Studierxs Lesx  Lesxs & x  xs (genitive) & xse & dix  einx & x \\
        * & Studier* Studier** Les*,  Les** & *  *s (genitive) & *'s & d*  ein* & * \\
        -ex & Studentex  Lesex & ex & ex & - & - \\
        -ens\footnotemark & Studentens Lesens & ens & ens & dens  einens & ens \\
\lspbottomrule
    \end{tabular}}
\end{table}


Some examples formulated with these strategies are:
\begin{itemize}
\setlength\itemsep{-0.5em}
    \item ``x sollte zukünftig besser besser auf xes Sprachhandlungen achten'' (they should in future be more careful with their choice of words).
    \item ``einx schlaux Studierx liebt xs Bücher'' (a smart student loves their books) (\cite[294]{hornscheidt2012feministische}).
    \item ``Lann liebt es, mit anderen zu diskutieren. Ex lädt häufig dazu ein, einen Roman zu besprechen'' (Lann loves to discuss things with others. They frequently invite people to discuss a novel).
    \item ``Lann und ex Freundex haben ex Rad bunt angestrichen'' (Lann and their friend have colourfully painted a bike) (\cite{hornscheidtex}).
    \item ``dens singend Radfahrens'' (the singing biker).
    \item ``dens solidarisch Rechtsberatens hat eine Kollektivausbildung absolviert'' (the solidary counsellor has completed a specialised education).
    \item ``einens hat das Recht zu gehorchen'' (they have the right to comply) (\cite[55]{hornscheidt2021wie}).
\end{itemize}
\footnotetext{Other possible forms are: \textit{keinens}, \textit{niemensch}, \textit{jemensch}, \textit{jedens}.}

The last approach to GFL in German is the development of neosystems that introduce a fourth gender in addition to masculine, feminine, and neuter, to make non-binary people visible in language and/or to refer to mixed groups of people. This approach consists of a new inventory of suffixes for (almost) each word class. A brief overview of some systems is given below.

The \textit{NoNa System} was (and is being) developed by two non-binary activists (\cite{NoNa}). As a personal pronoun, they propose the use of \textit{hen/hem}, derived from the Swedish gender-neutral pronoun \textit{hen} (see for example \cite{back2015hen}). Two sets of suffixes are introduced, i.e. \textit{-ai/-ais/-am/-ai} for definite articles, relative pronouns and demonstratives and\textit{ -t }plus a mix of usual masculine and feminine suffixes for indefinite articles and the negation word \textit{kein}, making it \textit{keint} (see Table \ref{tab:NonaSystem}). With respect to nouns, the activists suggest that the gender star is used to build them and the adjective declension mixes masculine and feminine suffixes (see Table \ref{tab:adjectivedeclensionnona} for an overview of the system with the phrase ``a good friend''). 

\begin{table}
    \caption[Pronouns and articles according to the NoNa System]{Pronouns and articles according to the NoNa-System (\cite{NoNa})}
    \label{tab:NonaSystem}
    \resizebox{\textwidth}{!}{
    \begin{tabular}{ccccccc}
    \lsptoprule
         \textbf{Case} & \textbf{Definite Articles} & \textbf{Indef. Articles} & \textbf{Personal Pronouns} & \textbf{Posses.} & \textbf{Demon.} & \textbf{Relatives}\\
\midrule
         \textbf{Nom.} & dai & eint & hen & meint & diesai & welchai \\
         \textbf{Gen.} & dais & einter & hens & meinter & diesais & - \\
         \textbf{Dat.} & dam & eintem & hem & meintem & diesam & welcham \\
         \textbf{Acc.} & dai & eint & hen & meint & diesai & welchai \\
\lspbottomrule
    \end{tabular}}
\end{table}

\begin{table}
        \caption[Declension system according to the NoNa System]{Declension system according to the NoNa System (\cite{NoNa})}
    \label{tab:adjectivedeclensionnona}
    \begin{tabularx}{\textwidth}{lQQ}
    \lsptoprule
         {Case} &  {Declension with indefinite   articles}  &  {Declension with definite   articles} \\
\midrule
         \textbf{Nom.} & eint gute Freund*in & dai gute Freund*in \\
         \textbf{Gen.} & einter guten Freund*in & dais guten Freund*in \\
         \textbf{Dat.} & eintem guten Freund*in & dam guten Freund*in \\
         \textbf{Acc.} & eint gute Freund*in & dai gute Freund*in \\
\lspbottomrule
    \end{tabularx}
\end{table}

Another system is the Sylvain system. During the Workshop \textit{Geschlechtergerechte Sprache} that took place as part of the \textit{Trans*Tagung} in Berlin in 2007, \textit{nin} was proposed as a gender-inclusive pronoun (see \cite{baumgartinger2008lieb}). This delivered the input for author De Sylvain to create a whole new grammar system to make German more inclusive and they also decided to write a novel using it \citep[40]{de2008sylvain}. This system introduced a new grammatical gender, the Indefinitivum (also liminales Geschlecht). Therefore, in addition to the terms \textit{der Mann} (the man) and \textit{die Frau} (the woman), \textit{din Lim} is to be used for people who reject the gender binary. An overview of the system can be found in Table \ref{tab:sylvain}.

\begin{table}[ht]
    \caption[Pronouns and articles according to the Sylvain system]{Pronouns and articles according to the Sylvain system (\cite{de2008sylvain})}
    \label{tab:sylvain}
    \resizebox{\textwidth}{!}{
    \begin{tabular}{ccccccc}
    \lsptoprule
        \textbf{Case} & \textbf{Def. Art.} & \textbf{Ind. Art.} & \textbf{Pers. Pron.} & \textbf{Poss. Pron.} & \textbf{Dem.} & \textbf{Rel. Pron.} \\
\midrule
         \textbf{Nom.} & din & einin & nin & meinin & diesin & welchin/din \\
         \textbf{Gen.} & dins & einins & nims/nimser & meinins & dienins & -/derin \\
         \textbf{Dat.} & dim & einim & nim & meinim & diesim & welchim/dim \\
         \textbf{Acc.} & din & einir & nin & meinin & diesin & welchin/din \\
\lspbottomrule
    \end{tabular}}
\end{table}

Other forms that are proposed include \textit{mensch}, \textit{jemensch}, \textit{niemensch}, and \textit{jedmensch} to replace the indefinite pronoun man,\footnote{The impersonal pronoun man because of its similarity to \textit{Mann} (i.e. man in English) is often regarded as a sexist element of the German language and some propose therefore to replace it.} \textit{jemand} (someone), \textit{niemand} (no one) and \textit{jedermann} (everybody). Finally, nouns are created by adding the suffix \textit{-nin} to the word stem in the singular (e.g. \textit{Studentnin}) and  \textit{-ninnen} for non-binary people only or\textit{ -Ninnen} for mixed-gender groups in the plural (e.g. \textit{StudentNinnen}). New adjectival suffixes for the indefinite article are proposed as well, as shown in Table \ref{tab:sylvaindeclension}, whereas feminine suffixes are used for the definite article. In the table, four phrases in German are reported, i.e. ``a young man'', ``a young woman'', ``a young non-binary person'', and ``a young animal''.

\begin{table}
        \caption[Declension system according to the Sylvain system]{Declension system according to the Sylvain system (\cite[45]{de2008sylvain})}
    \label{tab:sylvaindeclension}
    \resizebox{\textwidth}{!}{\begin{tabular}{ccccc}
    \lsptoprule
     & \multicolumn{4}{c}{\textbf{Singular}} \\
        \textbf{Case} & \textbf{Masculine} & \textbf{Feminine} & \textbf{Liminal} & \textbf{Neutral} \\
\midrule
        \textbf{Nom.} & ein junger Mann & eine junge Frau & einin jungin Lim & ein junger Tier \\
        \textbf{Gen.} & eines  jungen Mannes & einer jungen Frau & einins  jungen Lims & eines  jungen Tieres \\
        \textbf{Dat.} & einem  jungen Mann & einer jungen Frau & einim  jungen Lim & einem  jungen Tier \\
        \textbf{Acc.} & einen jungen Mann & eine junge Frau & einir jungin Lim & ein junges Tier \\
\lspbottomrule
    \end{tabular}}
\end{table}

\textcite{heger2020xier} introduced and has updated throughout the years the personal pronoun \textit{xier}. In 2009, the non-binary artist and activist started to engage with language to develop gender-inclusive alternatives to make German more inclusive for non-binary people (see \cite{hegerhistorie}). They proposed new neopronouns, constantly updating their forms. Currently, Heger uses \textit{xier}, derived from a fusion of the masculine and feminine pronouns \textit{er} and \textit{sie}. As shown in Table \ref{tab:xier}, definite articles, relative pronouns and possessives (as in ``their friend'', see the last column in the table) can also be built accordingly. 

%---------IMAGE begins here----------
%\begin{figure}[ht]
%\centering 
%\includegraphics[width=0.9\textwidth]{figures/Hegerpronomen.pdf}
%\caption[Gender-neutral pronouns proposed by Heger]{History of the development of gender-fair pronouns proposed by \textcite{hegerhistorie}}
%\label{fig:hegerpronouns}
%\end{figure}
%------------IMAGE ends here------------

\begin{table}
        \caption[Overview of the neopronoun \textit{xier}]{Overview of the pronoun \textit{xier} \citep{heger2020xier}}
    \label{tab:xier}
    \fittable{
    \begin{tabular}{llll}
    \lsptoprule
        \textbf{Case} & \textbf{Personal  Pronouns} & \textbf{Def. Articles/  Relative Pron.} & \textbf{Possessives}  \\
\midrule
        \textbf{Nom.} & xier & dier & xiesa Freund\_in \\
        \textbf{Gen.} & xieser & dies & xiesas Freund\_in \\
        \textbf{Dat.} & xiem & diem & xiesam Freund\_in \\
        \textbf{Acc.} & xien & dien & xiesan Freund\_in \\
\lspbottomrule
    \end{tabular}
    }
\end{table}

The debate on GFL also takes place on social media. For example, there is a private Facebook group called Geschlechtsneutrales Deutsch,\footnote{\url{https://www.facebook.com/groups/geschlechtsneutral/}} where about 500 members discuss language and try to find new solutions to overcome the binary structure of German.\footnote{Even though there is a third grammatical gender -- the neuter -- this is not used for people as it is seen as degrading and offensive.} Recently, the group administrators founded the Verein für Geschlechtsneutrales Deutsch (Association for gender-neutral German) \citep{vereindeutsch}. On its website, they propose three similar gender-fair systems, namely De-E-System, Dey-E-System and Dier-E-System. The first one is described to provide an overview of these more specific proposals. 

In the De-E-System, it is suggested that the gender-fair suffixes \textit{e} or \textit{re} as well as \textit{rne} should be used to build nouns respectively in the singular and plural. The personal pronoun should be \textit{en} \citep{vereindeutschdee}. Table \ref{tab:de-e-system} shows an overview of the declension of adjectives when they are preceded by no, definite and indefinite articles (within the phrase ``a good teacher''), and indefinite pronouns.

\begin{table}
    \caption[Overview of the \textit{De-E-System}]{Overview of the \textit{De-E-System} (\cite{vereindeutschdee})}
    \label{tab:de-e-system}
    \resizebox{\textwidth}{!}{\begin{tabular}{cccccc}
    \lsptoprule
      \multicolumn{1}{c}{\textbf{Case}} & {\textbf{Def. Art.}} & {\textbf{Indef. Art.}}& {\textbf{No Art.}} &\multicolumn{2}{c}{\textbf{Indef. Pronouns}} \\
\midrule
      \textbf{Nom.} & de  gute Lehrere & ein  gute Lehrere & gutey Lehrere & jedey & jemand  \\
      \textbf{Gen.} & ders  guten Lehreres & einers  guten Lehreres & guters Lehreres & jeders & jemanders \\
      \textbf{Dat.} & derm  guten Lehrere & einerm  guten Lehrere & guterm Lehrere & jederm & jemanderm \\
      \textbf{Acc.} & de  gute Lehrere & ein  gute Lehrere & gutey Lehrere & jedey & jemand \\
\lspbottomrule
    \end{tabular}}
\end{table}

As concerns general pronoun usage, in German pronouns \textit{er} and \textit{sie} are used to refer to people. There is no official equivalent to singular \textit{they} to refer to a person whose sex/gender is unknown (or irrelevant to the context of the conversation). Nevertheless, there are various solutions to this issue: 
\begin{itemize}
\setlength\itemsep{-0.5em}
    \item Impersonal, gender-neutral pronouns such as \textit{alle}, \textit{diejenigen}, \textit{jene} can be used (\cite[7]{unileipzig2020}).
    \item An individual's name can be repeated instead of using pronouns.
    \item Sentences can be reworded to avoid the use of pronouns.
    \item Gender-fair neopronouns can be used (see Table \ref{tab:genderfairpronouns} for an overview\footnote{In this case too, the overview is by no means complete, for other pronouns see \textcite{nibipronomen} and \textcite{pronomen}.}) -- note that trans*, non-binary, and inter people may use different pronouns. For this reason, one should always ask each person for their pronouns (\cite[38]{boka2021richtlinie}).
\end{itemize}

\begin{table}[!htb]
    \centering
    \caption[Gender-fair pronouns in German]{Gender-fair pronouns in German (\cite[36f]{hornscheidt2021wie}, \cite{heger2020xier}, \cite[294]{hornscheidt2012feministische}, \cite[33]{prona2016mein}, \cite[47]{de2008sylvain})}
    \label{tab:genderfairpronouns}
    \begin{tabular}{m{2cm}m{2cm}m{2cm}m{2cm}}
    \lsptoprule
         \textbf{Nom.} & \textbf{Gen.} & \textbf{Dat.} & \textbf{Acc.} \\
\midrule
         xier & xieser & xiem & xien \\
         er*sie & sein*ihr & ihm*ihr & ihn*sie \\
         y & y & y & y \\
         they & their & them & them \\
         dey & deren & demm & demm \\
         hen & hens & hem & hen \\
         ens & ens & ens & ens \\
         nin & nimser/nims & nim & nin \\
         x & xs & x & x \\
         ex & ex & ex & ex \\
         per & pers & per & per \\
\lspbottomrule
    \end{tabular}
\end{table}

Another field requiring GFL concerns forms of address. When addressing people, for example in emails, salutations are often gendered in German, e.g. ``Liebe Frau Smith'' not only refers to the person explicitly as \textit{Frau} or \textit{Herr}, but it also includes a gendered adjective since it is in the attributive position. Different strategies can be used to formulate a gender-fair form of address (\cite[77]{hornscheidt2021wie}): 
\begin{itemize}
\setlength\itemsep{-0.5em}
    \item Salutation followed by name and surname, such as \textit{hallo}, \textit{Sehr geehrt*}, \textit{Lieb*}, \textit{Guten Morgen/Tag}.
    \item Salutation only, e.g. \textit{Hallo}.
    \item \textit{Frau} and \textit{Herr} can be replaced by \textit{Person} or \textit{Mensch}.
    \item \textit{Guten Morgen} followed by full name.
    \item Gender star or other strategies can be applied, for example for job titles, e.g. \textit{Professor*in}.
\end{itemize}

The German language also comprises several terms which imply the gender binary. These include family terms, such as \textit{Mutter/Vater} (mother/father) or \textit{Tante/Onkel} (aunt/uncle) as well as job titles ending in \textit{-mann/-frau}, e.g. \textit{Feuerwehrmann} (firefighter) or \textit{Obfrau} (chairwoman). A list of gender-inclusive alternatives can be found on \textcite{nibispace}. Finally, a GFL is also a stereotype-free language, i.e. free from expressions suggesting how sex/genders are or should be, such as \textit{das starke Geschlecht} (the strong sex), that imply that men must be strong (\cite[33]{boka2021richtlinie}).

%--------------------------------------
%	SUBSECTION 3
%-------------------------------------
 
\subsection{Gender-fair language in Italian}
In Italian, the debate on GFL is still very much focused on the inclusion of women in discourse; more specifically, on the use of the feminine forms of functions and high-level professions from which women were once excluded. For instance, \textcite{gheno2019femminili} provides an overview of the rules to feminise functions and job titles in Italian and summarises as well as rebuts the arguments against them. The debate on language sexism was started by \textcite{sabatini1986raccomandazioni} who noted the widespread use of masculine generics, the asymmetrical use of name and titles, e.g. men called simply by surname whereas women by surname preceded by definite articles, e.g. ``la Tatcher'', and the lack of symmetry for offices and high-level professions. These are generally used in the masculine form even in reference to women.

Throughout the years, the debate on the inclusion of women in language has focused on this last point (see also the work of \cite{robustelli2012linee}, \cite*{robustelli2014donne}). Proposals to avoid the use of masculine generics have also been made, but they have been somewhat restricted to administrative texts; see for example the use of passive and impersonal\footnote{Nevertheless, gendering cannot always be avoided in these cases and some of the examples provided are illustrative: ``si entra uno alla volta'' -- i.e. in English ``enter one at a time'' -- one is a masculine form in Italian.} forms, indefinite pronouns, double forms, and gender-neutral nouns such as person \citep[22]{robustelli2012linee}.

More specifically, strategies to avoid masculine generics in Italian and to reach gender fairness are (\cite{vitiello2021linguaggio}):
\begin{itemize}
\setlength\itemsep{-0.5em}
    \item Change of subject so as to avoid past participles -- these are gendered in Italian.
    \item Use of periphrases.
    \item Use of synonyms for verbs, nouns, and adjectives.
    \item Change of perspective in the sentences.
    \item Use of neutral and collective nouns or impersonal pronouns, such as \textit{persona}, \textit{membro} (member), \textit{personale} (staff), \textit{chi} (who).
    \item Omission of nouns, pronouns and adjectives, and use of verbs as subjects.
    \item Use of passive and impersonal constructions (which are sometimes gendered, therefore one must be careful).
    \item Avoidance of the term \textit{uomo} (man) to refer to people in general. 
\end{itemize}
As mentioned in the section dedicated to the German language, such strategies can also be used to include non-binary people in discourse. 

With regard to the main topic of this work, i.e. the visibility of non-binary people, the asterisk was the first strategy introduced so as not to mark gendered word suffixes \citep{asterisk}. Curiously, it was introduced on the occasion of the Pride parade in Palermo in 2010 \citep{asterisk} and became quite popular on the internet \citep[50]{marotta2016linguaggio}, also gaining an entry in the grammar section of the Italian Encyclopedia \textit{Treccani} \citep{treccaniasterisco}. In expressions such as ``car* tutt*'' (dear all), the asterisk is used to avoid masculine generics and references to the gender binary \citep{treccaniasterisco}. Though this may represent a simple but not definitive solution for the written language, it cannot be orally transposed as it is not pronounceable \citep[50]{marotta2016linguaggio}. 

Other similar strategies have been proposed and are used (see Table \ref{tab:genderinclusiveitalian} for an overview). These are mainly graphical solutions, such as \textit{@} or \textit{-} and they are used to replace feminine and masculine suffixes in gendered elements like nouns, adjectives, and past participles as well as sometimes to build pronouns (see \cite{vitiello2021linguaggio}). 

\begin{table}
    \caption[Strategies for gender-inclusive Italian]{Strategies for gender-inclusive Italian based on \textcite{vitiello2021linguaggio}}
    \label{tab:genderinclusiveitalian}
    \begin{tabular}{lll}
    \lsptoprule
        \textbf{Ending} & \textbf{Example} & \textbf{Pronoun} \\
\midrule
        * & buongiorno a tutt* & l*i \\
        @ & buongiorno a tutt@ & l@i \\
        - & buongiorno a tutt- & \\
         & buongiorno a tutt & \\
        x & buongiorno a tuttx & lxi \\
        y & buongiorno a tutty & \\
        u & buongiorno a tuttu & \\
        ə/з & buongiorno a tuttз & ləi \\
\lspbottomrule
    \end{tabular}
\end{table}

Among all of the summarised strategies to represent non-binary individuals, the use of the International Phonetic Alphabet (IPA) symbol for the mid-central vowel schwa is worth considering. Proposed in 2015 by Boschetto \citep{itainclusivostoria}, an enthusiast on linguistics and inclusivity, it was made popular by \textcite[167]{gheno2019femminili}, who wondered whether such innovation would become the norm to refer to a multitude of people whose gender is unknown or irrelevant to the context of the conversation.\footnote{Note that \textcite[167]{gheno2019femminili} uses schwa (ə) to refer both to a single person and a multitude of people, e.g. ``carə collegə'' (dear colleagues) whereas Italiano Inclusivo (the site founded by Boschetto) suggests using another symbol (з) to maintain two distinct suffixes for singular and for plural \citep{itainclusivocome}.} Schwa has been adopted by publishing house effequ \citep{effequschwa} and has also been used in articles in the magazine \textit{L'Espresso} (see for example \cite{murgia2021espresso}) and in some books (e.g. \cite{meloni2021belle}\footnote{These authors do not follow suggestions by \textcite{itainclusivocome}; they rather use singular schwa as in \textcite[167]{gheno2019femminili}.}). Finally, it has been introduced among symbols in smartphone keyboards by Android and Apple \citep{itainclusivostrumenti}, but the latter adopted \textit{ə} only (see below). 

The use of schwa is very interesting for different reasons: (i) except for \textit{-u}, it is the only pronounceable strategy proposed; (ii) there is the possibility that singular and plural be maintained, as \textit{з} has been proposed for plural; (iii) it is a mid-vowel that is not generally present in the Italian language except for some dialects, hence it should not evoke maleness (for example the use of \textit{u}, which is a masculine ending in the Sardinian language) (\cite{itainclusivoperche}, \cite{gheno2021schwa}). In Table \ref{tab:schwa}, an overview of how to use schwa is proposed following suggestions by \textcite{itainclusivocome}. 

\begin{table}
    \caption[Overview of schwa in Italian]{Overview of schwa in Italian as proposed by \textcite{itainclusivocome}}
    \label{tab:schwa}
    \resizebox{\textwidth}{!}{\begin{tabular}{llll}
    \lsptoprule
        \textbf{Case} & \textbf{Feminine} & \textbf{Masculine} & \textbf{Inclusive}  \\
\midrule
        \textbf{Regular Nouns},  e.g. teacher & maestra & maestro & maestrə \\
        & maestre & maestri & maestrз \\
        \textbf{Irregular Nouns},  e.g. hero and painter & eroina & eroe & eroə \\
        & pittrici & pittori & pittorз \\
        \textbf{Epicene Nouns},  e.g. artist and poet & un'artista & un artista & un*artista \\
        & la poeta & il poeta & lə poeta \\
        \textbf{Articles} & la & il/lo & lə \\ 
        & le & i/gli & lз \\ 
        & una/un' & un/uno & unə/un* \\
        \textbf{3rd-Person  Singular Pronoun} & lei & lui & ləi \\
\lspbottomrule
    \end{tabular}}
\end{table}

Basically, schwa is used to replace feminine and masculine suffixes respectively in the singular and plural (as mentioned before, in Italian, adjectives, articles, pronouns, and past participles are also gendered). In the case of irregular nouns, that are hence not symmetrical, schwa replaces the last vowel of the masculine form. As regards articles, since the grammatical gender of the indeterminate \textit{un'} and \textit{un} can be distinguished by the apostrophe, the use of an asterisk (*) is suggested (\cite{itainclusivocome}). Nevertheless, a solution for gendered pairs such as brother/sister (in Italian there is no such term as sibling) or mother/father has yet to be proposed.

The use of schwa is also controversial (\cite{cosinonschwa,linguascema}, for example, strongly criticise it) and it is impossible to predict whether it will become the norm for GFL in Italian (see for example \cite{gheno2021schwa}). In particular, two main points are worth considering: schwa is considered to be ableist and it is used inconsistently across authors. Specifically, it is supposed to represent an obstacle for neurodivergent people reading it (\cite{gheno2021schwa}), though there are no sources or studies on the topic, and some screen readers cannot pronounce it \citep{orrù2021schwa}. Orrù, a copywriter and technical translator, explains other problems concerning the accessibility of schwa, its implementation in assistive technologies being the most prominent. Additionally, schwa has been appearing in books and articles (see the aforementioned \cite{effequschwa}, \cite{meloni2021belle}, \cite{murgia2021espresso}) but only the \textit{ə} has been used as in Table \ref{tab:alternativeschwa}. However, having two symbols, one for plural and one for singular, is part of what makes such a strategy appealing whereas the use of only one symbol may complicate reading. 

\begin{table}
    \caption[Alternative use of schwa]{Alternative use of schwa as in \textcite{effequschwanorme}}
    \label{tab:alternativeschwa}
    \begin{tabular}{lll}
    \lsptoprule
        & \textbf{Singular} & \textbf{Plural}  \\
\midrule
        \textbf{Nouns with articles} & lə maestrə &  ə maestrə \\
        \textbf{Nouns with prepositions} & dellə maestrə & deə maestrə \\
\lspbottomrule
    \end{tabular}
\end{table}

The present and the other two preceding subsections detailed specific strategies for GFL in English, German, and Italian. Note that GFL is language-dependent because grammatical structures differ among languages. Some good practices, however, may be applied regardless of the language and are briefly discussed in the following section. 
%--------------------------------------
%	SUBSECTION 4
%-------------------------------------
\subsection{Good practices in gender-fair language}
It is often thought that the sex/gender of a person can be deducted from external elements like physical appearance, voice, behaviour, clothing, and so forth. Based on that, people are referred to with masculine or feminine (pro)nouns \citep[32]{prona2016mein}. This may lead to misunderstandings that can be hurtful for people; research shows that being misgendered may cause feelings of stigma and psychological distress \citep{mclemore2018minority}. One should never assume to know people's sex/gender and should ask for identified pronouns instead. 

Another good practice consists in disclosing one's own identified pronouns, for example in the email signature or when introducing oneself to others \citep[23]{boka2021richtlinie}. Recently, social media such as LinkedIn and Instagram have introduced the option to add pronouns in the profile precisely in order to foster inclusivity \citep{linkedinpronouns,instagrampronouns}. This helps normalise the habit of asking for identified pronouns and acknowledges that many people's gender identity is not coincident with their sex assigned at birth. 

Finally, many expressions in language imply that a particular gender follows from one's sex assigned at birth. Expressions such as ``women with menstruation'' or ``male condoms'' imply that a person with menstruation identifies as a woman, whereas a person who uses external condoms is a man \citep[33]{boka2021richtlinie}. Such expressions can be avoided and replaced by neutral alternatives like ``people with menstruation'' or ``external condoms''. This shows how gender fairness encompasses knowledge of terminology concerning gender and sexuality as well.

\subsection{The societal impact of gender-fair language} \label{sec:impactgfl}

Subsection \ref{sub:malegenerics} addressed how the use of GFL can actually foster the visibility of women, especially in the workplace. The above-mentioned studies are fundamental in supporting the call for GFL. Nevertheless, such studies still concentrate either on binary linguistic strategies or only consider discrimination against women, thus ignoring non-binary individuals. Since the introduction and proliferation of the Swedish gender-neutral pronoun \textit{hen} (see \cite{gustafsson2015introducing}), some researchers focused on new, non-binary strategies in comparison to old neutral alternatives, finding that the latter still contain cognitive male bias.

\textcite{back2015hen} carried out two studies on the perception of \textit{hen}. In the first experiment, the researchers asked travellers at Stockholm Central Station to read the description of a person applying for a job as an insurance sales manager and to remember their gender. This profession was selected as it is equally occupied by men and women in Sweden. The description referred to the applicant as \textit{he}, \textit{she}, \textit{the applicant}, and finally \textit{hen}. Even though the Swedish equivalent for \textit{applicant} is also gender-neutral, it was found to evoke male imagery, whereas \textit{hen} was found to be far less biased. In the second experiment, other participants were asked to perform the same task but this time the job descriptions were presented using either typical masculine or feminine characteristics. The results were similar to those of the first experiment, and \textit{applicant} was found to be male-biased even when the job description was phrased with stereotypically feminine terms. 

\textcite{lindqvist2019reducing} conducted two other experiments, one in Swedish and the second in English, asking participants to read the description of a candidate for a job position. The candidate was referred to using (i) paired forms such as \textit{he/she} in English and in Swedish with the corresponding forms, (ii) traditional neutral words and singular \textit{they}, (iii) newly created gender-neutral forms such as gender-neutral pronoun \textit{hen} in Swedish and \textit{ze} in English. They found that in scenario (ii), people actually visualised men to a greater extent, hence traditional neutral wordings being cognitively biased towards men. Strategies (i) and (iii), instead, eliminated it. The scholars, however, call for the use of newly, actively created gender-neutral words instead of paired forms as the former are also inclusive of non-binary people. 

Similarly, \textcite{renstrom2023personal} tested whether paired and gender-neutral pronouns evoke a normative gender bias in Swedish and English. Normative bias is defined in the study as ``when words lead to stronger associations with individuals with normative gender expressions than with individuals with non-normative gender expressions'' \citep[2]{renstrom2023personal}. The first experiment was conducted in Swedish and participants were asked to read single sentences with the pronoun \textit{hen} or the paired pronoun \textit{he/she} used in reference to a person. The setup for the second experiment was similar and participants read about a candidate for a job position. Finally, a third experiment was conducted to replicate experiment two, but in English. The pronouns used were \textit{ze}, singular \textit{they}, and \textit{he/she}. The results suggest that neopronouns evoke fewer normative associations when they are newly created and recognized by language users as non-binary pronouns, since they are not linked to masculine or feminine genders..

Some initial studies on the effects of German gender star on mental representation have appeared, e.g. \textcite{kurz2023star} who replicated the study by \textcite{stahlberg2001effekte} presented in Subsection \ref{sub:malegenerics}. The researchers focused on the use of masculine generics and gender star and found that this last strategy leads to more inclusive mental representations.

To the best of my knowledge, these are some of the few studies in which gender-neutral, non-binary strategies are taken into account, although they do not always focus on the effects of such strategies on the visibility of non-binary individuals. It can only be presumed that their usage will help dismantle the gender binary in the long run, creating a more respectful (not only) linguistic environment for non-binary people.


\section{Research on gender-fair language in (machine) translation} \label{relatedwork}
GFL poses a challenge for human and MT alike. Since non-binary people are featured in different typologies of popular texts and formats, ranging from magazine articles to television series and films, translators must be able to opt for appropriate strategies that reflect the language used in the source text \citep{misiek2020misgendered} while translating and/or post-editing. In addition, MT systems still suffer from gender bias, thus gender identities beyond the binary are underrepresented or erased in MT outputs. Issues of the representation of non-binary people in language are neglected in both TS and MT. 

This section summarises selected work focusing on non-binary GFL either from a TS, an MT, or both perspectives. For a more comprehensive literature review, see \textcite{nettt2022}. Note that GFL research is still in its infancy, both in TS and MT, and numerous studies are appearing at the time being, hence the overview provided in this section is by no means exhaustive. 

\subsection{Gender-fair language in translation studies} \label{sec:genderfairtranslation}

\textcite{misiek2020misgendered} analyses the Polish translation of three different English-language TV series that feature non-binary characters. 
The scholar finds that translators often misgendered the characters or even completely erased their gender identity. Polish is a grammatical gender language which has the feminine, masculine, and neuter. It has three third-person singular pronouns, i.e. \textit{on}, \textit{ona}, and \textit{ono}, which respectively correspond to English he, she, and it. It has also two third-person plural pronouns, i.e. \textit{oni} for the masculine and one for the feminine and neuter. There is hence no direct equivalent to singular \textit{they} to refer to non-binary individuals. Misiek also provides an overview of the strategies for gender neutrality in Polish. They include distortion and/or omission of verb suffixes, use of passive constructions, present tense because the past is gendered, use of the vowel \textit{u} as a new neuter suffix, or use of neuter forms \citep[167]{misiek2020misgendered}. As regards translation, the scholar proposes to: (i) always ask the actor interpreting the character how they are to be referred to in Polish; (ii) consult non-binary or genderqueer communities; (iii) use neuter forms or the aforementioned vowel \textit{u}; (iv) use neologisms to describe professions; and (v) use general, non-gendered words like person.

\textcite{vsincek2020ona} focuses on Croatian and analyses a film translation as well as journal articles on Sam Smith's coming out as non-binary. The researcher also finds that non-binary gender identities are misrepresented or even erased in both the film and the news reports. Similar to Polish, Croatian is a grammatical gender language with three third-person singular pronouns, i.e. \textit{on}, \textit{ona}, and \textit{ono} which correspond to English \textit{he}, \textit{she}, and \textit{it}. In contrast to Polish, there are three plural forms: oni is masculine, one feminine, and ona neuter, hence there being no equivalent to singular \textit{they} in Croatian either. When analysing Sam Smith's coming out as non-binary, Šincek found that they were referred to as \textit{on}, \textit{ono} or \textit{oni} -- this last being a literal translation of singular \textit{they}. The researcher also conducts interviews with three non-binary people asking which pronouns they use in Croatian. One participant uses a pronoun that reflects their sex assigned at birth but thinks that \textit{ono} could be a good gender-neutral pronoun whereas the other two consider it derogatory because it is used to describe inanimate objects. Another participant tried to use \textit{oni} for some time but noted how this led to confusion. In general, participants agree to use archaic verb forms that are not gendered in Croatian as well as gender-neutral words such as person. To date, there are no dedicated strategies to achieve gender fairness in Croatian, and English gender-neutral pronoun singular \textit{they} is often rendered as oni, which is nevertheless third-person plural and masculine. Šincek proposes (i) to use simple tenses since these are not gendered; (ii) as well as the word person; and (iii) to try to add a new meaning to oni. 

López's (\cite*{lopez2019tu,lopez2022trans}) work focuses on the translation of media products as well. These include the Spanish dubbed and the subtitled versions of English language TV shows, such as \textit{One Day at a Time}, which feature non-binary characters. López distinguishes between INL and DNL. In the former, sentences are structured in a way that gender can be avoided, e.g. using passive constructions and gender-neutral terms such as person. INL may be used to address a broad audience that includes different genders since it conceals them. In the latter, new gendered suffixes, typographical characters, and neopronouns may be used to specifically address non-binary individuals. The researcher finally notes that (i) the identity of non-binary characters is often erased in the TV shows analysed; (ii) gendered forms are used inconsistently; and (iii) non-binary language is often translated with unusual forms. Similarly, \textcite{attig2023call} analyses the Spanish and French dubbed and subtitled versions of \textit{One Day At a Time} and calls for a community-informed translation practice of queer-oriented texts.

\textcite{lardellitranslating} provide an overview of GFL strategies beyond the binary in German and Italian. To investigate their actual use, they create a corpus with online reports about Demi Lovato's coming out as non-binary in both languages. They find that misgendering is still very common, especially in Italian language texts. Translation solutions for singular \textit{they} greatly differ and show that journalists often lack knowledge about non-binary genders and language. While a dominant GFL strategy was not found in the analysed corpus, a mix of gender-neutral and gender-inclusive forms was used.

The present work employs a process-oriented framework to study the use of GFL in translation. Specifically, it analyses the impact of different strategies on cognitive processing, focusing on times, screen activity tracking data, and retrospective accounts. Such study is important because it contributes towards bringing translation professionals to the forefront of the discourse around gender fairness in language and translation, also showcasing their role as mediators and agents for social change. The participants in the study were not regarded as passive agents, but they were asked to experiment with different strategies to gain insights into their professional and personal opinions regarding GFL. 


\subsection{Gender-fair language in machine translation}
\textcite{saunders2020neural} propose a first approach to address gender-fair MT beyond the binary. The researchers extended a gender-balanced dataset (\cite{saunders2020reducing}) with gender-neutral sentences containing placeholders for gender inflexions in German and French. They trained an NMT system through transfer learning and tested the effectiveness of their methods with a gender-neutral version of WinoMT (\cite{stanovsky2019evaluating}) that they produced. The NMT model scored a low overall accuracy, struggling with the generation of gender-neutral language.

\textcite{lauscher-etal-2023-em} analyse how three commercial MT systems translate third-person pronouns. They compare the translations of gendered, gender-neutral, and neopronouns from English to Danish, Farsi, French, German, and Italian as well as from Danish to English. They find that gender-neutral pronouns often lead to grammatical and semantic translation errors, e.g. pronouns are often wrongly translated, omitted, or even copied in the MT outputs. In addition, gender neutrality is often not preserved leading to misgendering.


In the literature, there is a clear trend toward developing diverse benchmarks, spanning multiple languages and focusing on either gender-neutral, gender-inclusive, or non-binary language \citep{piergentili-etal-2023-hi,piergentili-etal-2024-enhancing,lardelli-etal-2024-building}. Moreover, the field is increasingly focusing on LLMs, for instance testing their capabilities to produce GFL while bias evaluation methods are evolving from traditional reference-based scoring to LLM-driven, reference-free approaches, often leveraging chain-of-thought prompting for improved accuracy \citep{piergentili2025llm}. A brief overview of such research directions is provided in the following paragraphs.

\textcite{piergentili2023gender} tackle the issue of gender bias in MT from a theoretical perspective. While they do not specifically address non-binary language issues, they focus on gender-neutral language strategies such as indefinite pronouns, omissions, repetitions, and gender-neutral alternatives to gendered terms. These strategies should be used, for example, to avoid assigning a gender in the target text when the source is ambiguous, and to avoid the generic masculine. The scholars reviewed gender-neutral strategies in English and Italian, outlined a definition of gender-neutral translation and discussed the technical challenges involved in implementing it in MT such as the dynamicity of the gender-neutral strategies, the lack of training data and methods as well as evaluation procedures. Based on this, \textcite{piergentili-etal-2023-hi} introduce GeNTE, a benchmark for evaluating gender-neutral translation from English to Italian. They propose both reference-based and reference-free evaluation methods and show that current MT systems still struggle in this area, whereas \textcite{piergentili-etal-2024-enhancing} explore gender-inclusive translation into Italian using neomorphemes (e.g. \textit{tuttə}) and LLMs. They release NEO-GATE, a benchmark for evaluating gender-inclusive translation, and demonstrate that LLM prompting can improve inclusive language generation. Finally, \textcite{piergentili2025llm} propose an LLM-as-a-judge method for evaluating gender-neutral translations in a scalable, reference-free manner. Using phrase-level prompting improves LLM judgment reliability across multiple languages and enhances automation of gender bias evaluation in MT.

\textcite{lardelli-etal-2024-building} present a benchmark for gender-fair translation from English to German using real-world texts. Their analysis reveals a strong tendency toward masculine defaults and limited use of inclusive language, even among state-of-the-art systems. \textcite{robinson-etal-2024-mittens} introduce MiTTenS, a multilingual benchmark covering 26 languages to evaluate gender mistranslation. The dataset includes both naturalistic and synthetic data targeting known failure modes, revealing persistent gender bias across MT systems and language pairs.

Finally, \textcite{chen2024beyond} and \textcite{jourdan2025fairtranslate} strongly focus on gender bias against non-binary people. The former propose AmbGIMT, a multilingual benchmark measuring gender bias in translation using ambiguous attitude expressions. They introduce the Emotional Attitude Score (EAS) and find that translations of non-binary contexts often carry more negative sentiment and reduced quality. The latter release FairTranslate, a gender bias evaluation dataset for English–French MT with a focus on occupational terms and inclusive forms. Through LLM prompting experiments, they show that non-binary forms receive lower translation quality and exhibit persistent stereotypes.

\textcite{amrhein-etal-2023-exploiting} propose a gender-fair rewriting model in line with other research in the field of debiasing MT (see, e.g. \cite{vanmassenhove-etal-2021-neutral}). While in previous studies training data for these models were created through linguistic rules, \textcite{amrhein-etal-2023-exploiting} use MT models to create gender-biased text from real gender-fair text via round-trip translation. They train a rewriting model for German which was assessed in a human evaluation study and showed an increase in gender fairness in comparison with biased MT outputs. 

The present study differs from the research presented in this subsection for various reasons. It neither analyses the outputs of various commercial MT systems nor proposes debiasing approaches. Instead, it examines the PE process with the aim, among others, of determining whether PE is a viable instrument for producing gender-fair datasets that could further research on gender-fair MT. Additionally, it collects the opinions of language professionals to call for more user-centred and interdisciplinary-informed debiasing approaches, potentially prompting a shift in how gender bias research is conducted within computational linguistics.


\subsection{Interdisciplinary research on gender-fair language}
\textcite{burtscher-2022-genderfairmt} and \textcite{gromann2023participatory} propose a method building on participatory action research to include experiential experts -- i.e. queer and non-binary people, translators, and MT experts -- in the discussion around gender fairness in human translation and its implementation in the MT design process. The researchers organised a three-day participatory workshop focusing on (machine) translation from English into German. Some of the major findings include the importance of context dependency to avoid identity invalidation, the need for materials and resources on GFL, and a desire for customisable MT solutions.

The present work contributes to discussion among TS and MT researchers by investigating both the translation and PE process of non-binary GFL. Similar to existing research, the focus is on the translation from English into languages for which the debate on GFL is much more developed. Although to a different extent and with varying degrees of success, in German-speaking countries GFL and GFL policies have been discussed since the 1970s (see \cite{sczesny2016can}), and several strategies to represent and/or include non-binary individuals in discourse have been developed, as explained in Section \ref{inclusivelanguage}. 

While related research in TS generally provides overviews of common strategies for GFL and analyses their use in the translation of different media products, the present study combines process and product data to provide more comprehensive insights into the translation and PE process which can support translation professionals in their work and also provide inspiration for the debiasing of MT. 

\section{Relevance for the empirical study}
The scope of this final section is to summarise key theoretical aspects, highlighting their relevance for the empirical study and analysis proposed in this book.

An overview of various concurrent GFL strategies was provided, illustrating different approaches to gender fairness both within and across languages. Languages have different grammatical structures, thus complicating the translation of gender, as also presented in Subsection \ref{sub:translatinggender}. A source language may express gender neutrally or ambiguously, whereas the target language may not. In such cases, GFL may be needed, but the choice of strategy can meaningfully affect the translation’s interpretation.

A fundamental principle in the translation practice and TS is the existence of multiple translations for a single source text. Different versions may be generated based on the translation's intended purpose \citep{reiss1984grundlegung}. Furthermore, translation can be seen as a process of intercultural communication, resulting in a text that functions appropriately in specific situations and contexts \citep{holz1984translatorisches}. 

Consequently, participants were instructed to experiment with three GFL approaches to facilitate comparative analysis and discussion on the implications of their usage. During the retrospective interviews, the concept of the context-dependent selection of strategies was discussed, with participants elucidating the impact of GFL on the target texts (see Chapter \ref{Discussion}). 

Additionally, the literature review presented in this chapter revealed a lack of interdisciplinarity, with researchers from TS and computational linguistics often overlooking studies outside their respective fields. To date, the only interdisciplinary research efforts have been collaborative projects (\cite{burtscher-2022-genderfairmt,gromann2023participatory}). This study aimed to investigate the perspectives of translation professionals and integrate viewpoints from TS, a discipline that has historically overlooked research on translation technology (\cite{rozmyslowicz2024politics}).

Finally, excerpts from this chapter were used to create a handout on GFL, which was distributed to study participants to prepare them for the GFL translation and PE tasks. These contents also served as a basis for analysing the strategies employed and the errors made by participants, as discussed in Chapter \ref{methodology}.
